\section{\large\textbf{Introduction}}

The use of surround sound systems in non-anechoic environments can alter and impair the user experience. Except for rooms specifically treated for this purpose, the room's inherent acoustic response, whether due to unwanted specular reflections or resonant modes, distorts the information received by the listener, causing erroneous perception of source locations. The use of parameters such as the directivity of the loudspeakers, allows to reduce the energy contribution of the room, increasing the direct-to-reverberant energy ratio (\textit{D/R}), minimizing their effects. However, achieving such directivity control becomes difficult as frequency decreases.

Sound source localization has been a subject of study for a long time. \textcite{Rayleigh} examined the implications of interaural time difference (\textit{ITD}) and interaural level difference (\textit{ILD}) parameters in humans' ability to localize sound sources. These parameters for distant sources exhibit interaural ambiguity, leading the author to define a "cone of confusion". Furthermore, this study has approximated the frequency ranges where these parameters are useful; at high frequencies (above 2 kHz), an acoustic shadow is created by the head, generating interaural level differences, while at mid-low frequencies (below 1 kHz), this level difference is minimal, but an interaural time difference is generated between the two ears. More recent studies have determined that the perceptible high-frequency limit for \textit{ITD} is 1.4 kHz, while below 700 Hz it corresponds to a constant interaural phase difference, given that the wavelength is significantly greater than the head's diameter \parencite{hflimit}. Additionally, \textcite{mills} has studied these parameters and determined minimum perceptible differences of 10 $\mu$s for \textit{ITD} and 0.5 dB for \textit{ILD}, under optimal conditions. The type of signal used is critical for localization ability; the duplex theory highlights the predominance of \textit{ITD} over \textit{ILD} in broadband signals, thus making low-frequency information dominant \parencite{duplex}. However, since level differences can also arise from the distance difference between each ear and the source, this parameter prevails over \textit{ITD} for sources located less than two meters away, where in these cases, a torus of confusion centered on the interaural axis is defined \parencite{tori}.

The previously cited studies were performed using anechoic stimuli in order to avoid any interference that might alter the perception of these phenomena in controlled conditions. However, the typical listening environment tends to be reverberant, which means the effect of the room is underestimated, whereas in reality, it significantly impacts perception \parencite{rev1}. Due to the temporal alterations that modify the transmitted sound, it is crucial to understand the envelope of the signal to be localized; \textcite{Hartmann} conducted studies using impulsive and non-impulsive signals in a variable acoustic room. For impulsive signals, the room’s geometry was found to be decisive and independent of the reverberation time, while for non-impulsive signals, reverberation time and the \textit{D/R} ratio were determinant factors. Additionally, broadband signals were the easiest to localize. On the other hand, \textcite{espectrorev} concluded that \textit{ITD} predominates over \textit{ILD} in reverberant spaces, even though level differences provide more precise information than temporal differences.

The aim of this article is to study the localization of sound sources in a reverberant field in the low-frequency range (below 1.4 kHz), modifying the directivity of the source and, therefore, the \textit{D/R} ratio. Unlike previous studies, the listener will always be far from the source, thus avoiding a considerable change in \textit{ILD} that is not related to the source-room interaction. This study seeks to provide useful information for the design of home electroacoustic systems to optimize both frequency and directivity responses, based on the intended use of the equipment.
